% 导读：译者所加，不是原文内容
\section*{导读：这篇文章在讲一个什么故事}
\addcontentsline{toc}{section}{导读：这篇文章在讲一个什么故事（译者所加）}
\markboth{导读}{导读}

\begin{keytakeaways}
\item 一句话：本文证明，在「A 模型与 B 模型互换」的视角下，射影丛
      $\mathbb{P}(\mathcal{O}_X\oplus L)$ 的镜像不是另一个射影丛，而是一个\textbf{带停止位的映射柱}
      $M(X,\phi)=X\times_\phi\mathbb{C}^*$（Kartal 修改的 Weinstein 结构）上的包裹 Fukaya 范畴。
\item 这条定理把「同调镜像对称（HMS）」从「一对流形」推广到「一族退化的流形」，
      并顺带造出一族两两不等价的异域 Weinstein 结构。
\item 读法建议：只想抓故事，读 §1 与本导读即可；想懂证明，按 §2 → §4 → §5 → §6 的顺序读；
      附录 A 是 §6 的技术工具，可最后看。
\end{keytakeaways}

\subsection*{一、出发点：HMS 说了什么，又没说什么}

同调镜像对称（homological mirror symmetry, HMS）是 Kontsevich 1994 年提出的猜测：对一对
互为镜像的辛流形 $(X,\omega)$ 与代数簇 $Y$，辛流形一侧的 Fukaya 范畴与代数簇一侧的凝聚层
导出范畴应当等价，
\[
  \mathcal{F}(X,\omega)\;\simeq\;\Coh(Y).
\]
左边（A 模型）只记得辛结构：对象是 Lagrangian 子流形连同局部系统，态射由 Floer 同调给出；
右边（B 模型）只记得复结构：对象是凝聚层复形，态射是 Ext。这个猜测的威力在于，它把
「几何上互为镜像」这句模糊的话翻译成了精确的范畴等价，从而让两边的计算可以互相搬运。

二十多年过去，HMS 已经在椭圆曲线、四次曲面（Seidel）、环簇 Fano 流形与对数 Calabi–Yau
曲面（Auroux–Katzarkov–Orlov 等）等情形得到验证。但这些验证大多是「单个流形对单个流形」的。
本文要问的是\textbf{族的问题}：当 $X$ 退化成一族 $X_t$ 时，镜像那一侧应当发生什么？

\subsection*{二、卡点：环簇退化的镜像问题}

设 $(X,Y)$ 是一对镜像，$X$ 有一个环簇退化 $X\rightsquigarrow X_0$（例如把 $\mathbb{P}^1$ 收成
两个点、把曲面收成若干环面块的并）。SYZ 纲领告诉我们：$X$ 的镜像 $Y$ 上有一个
Lagrangian 环面纤维化，$X$ 的退化对应 $Y$ 上的一个「热带」极限 $\omega_p(t)$，其中参数 $t$
记录退化程度。自然的猜测是
\[
  \Coh(X_t)\;\simeq\;\mathcal{F}\bigl(Y,\omega_p(t)\bigr)\qquad(\text{对所有 }t),
\]
即「退化只改变辛形式、不改变镜像的拓扑」。这句话听起来合理，却极难直接证明：$t\to 0$ 时
辛形式趋向退化，Fukaya 范畴的 Floer 理论在极限处失效。

本文的实际切入点是\textbf{对数几何}：把退化看成沿边界的对数结构，对数爆破（log blow-up）是
在边界附近做手术的标准操作。如果知道对数爆破的镜像是什么，就能把 $X_t$ 的镜像一步步
算出来。而对数爆破的局部模型是射影丛 $\mathbb{P}(\mathcal{O}_X\oplus L)$，于是问题归结为：

\begin{quote}
\bfseries 对一个 HMS 对 $(X,Y)$ 与线丛 $L$，$\mathbb{P}(\mathcal{O}_X\oplus L)$ 的镜像是什么？
\end{quote}

这正是原文问题 1.3，也是全文的主线。注意这个问题本身是「局部」的：它只关心一个射影丛，
而不关心整个退化族；一旦能对任意线丛回答它，就能靠粘合（扇形下降）把退化的各块拼起来。

\subsection*{三、答案的形状：为什么是映射柱}

在一维情形先找直觉。$\mathbb{P}^1$ 的镜像是 $(\mathbb{C}^*,z+\frac1z)$：一个柱面，上面有一个
「超势」$z+\frac1z$，其临界点给出两个停止位（stop），分别对应 $\mathbb{P}^1$ 的 $0$ 与 $\infty$。
更一般地，线丛 $L$ 的「绕一圈」信息由 $\mathbb{C}^*$ 作用给出：把 $X\times\mathbb{C}^*$ 沿
$\phi$（$L$ 对应的镜像自同构）粘起来，就得到映射柱
\[
  M(X,\phi)\;=\;X\times_\phi\mathbb{C}^*,
\]
它以 $\mathbb{C}^*$ 纤维化于 $S^1$，绕 $S^1$ 一周的单值性恰是 $\phi$。再在 $M(X,\phi)$ 上取
「两条纤维」作为停止位（由 $z+\frac1z=c$ 给出），得到部分包裹 Fukaya 范畴
$\mathcal{W}(M(X,\phi),z+\frac1z)$。

本文主定理（定理 1.4）说：在 HMS 对 $(X,Y)$ 与镜像自等价 $(\phi,\xi)$ 的假设下，存在
$A_\infty$-函子
\[
  \Coh\bigl(\mathbb{P}(Y,\xi)\bigr)_{\mathrm{dg}}
  \;\simeq\;
  \Perf\,\mathcal{W}\Bigl(M(X,\phi),\,z+\tfrac1z\Bigr)^{\mathrm{op}},
\]
它在同伦范畴上诱导三角范畴的等价。这里 $\mathbb{P}(Y,\xi)$ 是 $Y$ 上由 $\xi$ 扭的射影丛
（B 侧），$M(X,\phi)$ 是 A 侧的映射柱。等式的形状值得玩味：\textbf{射影丛的方向被「翻译」成了
映射柱方向}——$\mathbb{P}^1$-丛的 B 模型，对应一个带单值性的 A 模型。

\subsection*{四、推论：从线丛一路到 Hirzebruch 面}

把主定理喂给具体的 $X$，就得到一串推论（原文 §1b–§1c）：

\begin{itemize}
\item \textbf{线丛与主 $\mathbb{G}_m$-丛}（推论 1.5）：$L$ 的镜像自同构 $\phi$ 直接给出
      $\Coh(L)$ 与 $\Perf\mathcal{W}(M(X,\phi))$ 的等价；取 $L=\mathcal{O}$ 回到原来的 HMS。
\item \textbf{多重单值性}（推论 1.6–1.7）：把 $k$ 个 $\mathbb{P}^1$ 沿节点粘成一条链（nodal chain），
      镜像不再是单个映射柱，而是一串映射柱的粘合；用扇形下降与真下降把各块拼起来，
      得到 $\Coh$ 侧的「反常 Schober」与 A 侧的粘合范畴之间的等价。
\item \textbf{环面丛}（推论 1.8）：$\mathbb{P}(\mathcal{O}\oplus L_1\oplus\cdots\oplus L_r)$ 的镜像
      是多重映射柱，对应环面丛的 A 模型。
\item \textbf{Hirzebruch 面}（推论 1.9）：$F_n=\mathbb{P}(\mathcal{O}_{\mathbb{P}^1}\oplus\mathcal{O}(n))$
      的镜像是 $\mathbb{P}^1$ 上某个带停止位的映射柱；由于 $\phi$ 的「扭转量」随 $n$ 变化，
      这些 A 模型两两不等价。
\item \textbf{异域 Weinstein 结构}（推论 1.10）：$F_n$（$n\ge0$）作为辛流形两两不精确辛同构，
      但其包裹 Fukaya 范畴却互相等价——这说明 A 模型看不见 B 侧的复结构差异，也给出
      「同一光滑流形上无穷多个互不等价的 Weinstein 结构」的显式例子。
\end{itemize}

\subsection*{五、技术骨架：两边各需要什么}

\textbf{A 侧（辛几何）}的核心新工具是「离散商的包裹 Fukaya 范畴」。映射柱 $M(X,\phi)$ 是
$X\times\mathbb{C}^*$ 对 $\mathbb{Z}$（或更一般的离散群）作用的商，但 $\mathcal{W}$ 不能直接
在商上定义（Reeb 流不下降）。作者的办法是先在覆盖上定义带群作用的 Weinstein 扇形上的
Fukaya 范畴，证明它有群作用、再取轨道范畴，从而得到 $\mathcal{W}(X/G)\simeq\mathcal{W}(X)/G$
型的等价（定理 4.1）。此外还需要：商的 Künneth 公式（定理 4.4）、Kartal 修改的
「$z+1/z$」结构（§4d），以及扇形下降/真下降的粘合机器（§1b、§2b）。

\textbf{B 侧（代数几何）}的核心新工具是「$G$-等变层的导出范畴」。$\mathbb{P}(\mathcal{O}\oplus L)$
带一个 $\mathbb{G}_m$ 作用（旋转纤维），其凝聚层按权分解；作者系统建立了等变层的内 Hom、
权空间分解、限制到子群概形、群作用的乘积等工具（§5），最后证明扭 Künneth 公式（§6c）：
\[
  \text{等变层的导出范畴}\;\simeq\;\text{（带扭的）普通导出范畴}.
\]
把 B 侧的这套语言与 A 侧的商 Fukaya 范畴对齐，就是主定理的证明。

\textbf{附录 A} 用 Eilenberg–Zilber 收缩为 Čech 复形提供同伦共轭，是 §6 里「扭」的
技术基础；它本身是古典结果，读起来最「干」，可以最后看。

\subsection*{六、证明路线图}

原文 §2 是全文中枢，它把主定理拆成三步：
\begin{enumerate}
\item 先用扇形下降把 $\mathbb{P}(\mathcal{O}\oplus L)$ 分解为若干「扇形块」上的问题
      （§2b，定理 2.1–2.2）；
\item 对每个块，用环面丛的情形（§2c）算出镜像；
\item 最后用 Hirzebruch 面 $F_n$ 的具体计算（§2d）验证拼接与单值性的一致性。
\end{enumerate}
可以说，§1 陈述「故事」，§2 证明「故事」，§3–§6 则是为 §2 里每一步准备的语言。

\subsection*{七、怎么读这篇文章}

\begin{itemize}
\item \textbf{只想知道结论与故事}：读 §1 与本导读，再扫一眼 §2 的定理陈述即可。
\item \textbf{想理解主定理}：先读 §4（A 侧商的 Fukaya 范畴）与 §5（B 侧等变层），
      把两边的「语言」各过一遍，再回头看 §2 的证明。
\item \textbf{想做计算}：直接看 §2d（Hirzebruch 面）与 §1c 的例子，那里公式最具体。
\item \textbf{关心技术细节}：§3（dg/A∞ 与轨道范畴）与 §6（扭 Künneth）是工具章，
      可与 §4、§5 交叉阅读；附录 A 备查。
\item \textbf{术语}：所有专门术语在文末「术语表」有中英对照；首次出现的术语在正文括注英文。
\end{itemize}

\subsection*{八、本文在文献中的位置}

本文站在几条线索的交汇处：Kontsevich 的 HMS 与 SYZ 纲领给出了问题；FLTZ（Fukaya–Oh–Ohta–
Ono 型的环簇镜像）与 Gross–Siebert 的对数/热带纲领提供了「退化与粘合」的语汇；
Abouzaid 与 Ganatra–Pardon–Shende 的包裹 Fukaya 范畴理论提供了 A 侧工具；
Kartal 的映射柱计算给出了关键技术输入；Sylvan 的部分包裹范畴与 Kapranov–Schechtman 的
反常 Schober 则分别对应 A、B 两侧的粘合语言。本文的贡献是把这些线索在「射影丛的镜像」
这一个点上拧成定理，并由此得到异域 Weinstein 结构这一副产品。

\subsection*{九、留白}

本文的主定理要求 $(X,Y)$ 已经是 HMS 对、且 $(\phi,\xi)$ 是镜像自等价——也就是说，它
\textbf{不构造}新的镜像对，而是「在已知镜像对上做手术」。对一般的环簇退化 $X_t$，还需要
验证各块的镜像数据能拼成全局的 $\phi$；原文把这一步留作后续（问题 1.11 及 §1d 的讨论）。
此外，$\mathbb{P}^1$ 换成高维射影空间、或 $L$ 换成秩 $>1$ 的向量丛，本文只给出部分情形
（§1c 的环面丛），一般情形仍是开放的。
